\documentclass[11pt]{article}

\usepackage[a4paper,margin=2.5cm]{geometry}
\usepackage{amsmath,amssymb}
\usepackage{graphicx}
\usepackage{hyperref}
\usepackage{url}
\usepackage{enumitem}

\title{GRA Forum: Orchestrating AI Agent Debate via Gradient Reduction of Argumentative Foam}

\author{AAA AAA%
\thanks{ORCID: \href{https://orcid.org/0009-0004-1872-1153}{0009-0004-1872-1153}.%
Repository: \href{https://github.com/qqewq/gra-forum}{github.com/qqewq/gra-forum}.%
Zenodo: \href{https://doi.org/10.5281/zenodo.19158686}{10.5281/zenodo.19158686}.}}

\date{March 2026}

\begin{document}

\maketitle

\begin{abstract}
Large language models (LLMs) and other AI agents can be combined into multi-agent systems that debate, collaborate and critique each other.
However, naive orchestration of such systems quickly degenerates into redundant chatter, unresolved contradictions and unstructured noise.
This paper introduces \emph{GRA Forum}, a minimal, practical orchestrator for AI agent debate based on the principle of \emph{Gradient Reduction of Argumentative foam} (GRA).
The core idea is to treat multi-round debates as a dynamical system with explicitly measured ``foam''---a set of scalar functionals capturing conflict, vacuity, redundancy and noise---and to iteratively steer the debate by minimizing a global objective \(J\) over these metrics while preserving constructive disagreement.
The implementation is released as an open-source Python library.\footnote{\url{https://github.com/qqewq/gra-forum}}[page:0]
\end{abstract}

\section{Introduction}

Recent progress in large language models and tool-augmented agents has enabled complex multi-agent systems where different models or roles interact in natural language.
Such systems are used for self-consistency, tool-use orchestration, code review, scientific reasoning and more.
Yet most existing multi-agent frameworks focus on wiring agents together (who talks to whom) rather than on \emph{controlling the dynamics of argumentation} itself.[page:0]

In practice, unregulated agent debates exhibit several pathological phenomena:
(i) unresolved contradictions between agents,
(ii) vacuous or ungrounded statements,
(iii) redundant restatements of the same content, and
(iv) low-level noise such as formatting errors or inconsistent output schemas.
We refer to the aggregate of these phenomena as \emph{argumentative foam}.[page:0]

The GRA Forum project proposes a simple but general principle: treat argumentative foam as a measurable quantity \(\Phi\) on the evolving debate state, define a global cost \(J\) over several foam components, and implement an orchestrator that adapts prompts, questions and agent routing to reduce parasitic foam over successive rounds.[page:0]
Unlike approaches that aim for immediate consensus, GRA Forum explicitly preserves constructive disagreement as a source of epistemic gain, while penalizing only non-productive components of foam.[page:0]

\section{System Overview}

GRA Forum is organized as a three-level architecture (with possible extensions to additional levels such as physics-informed agents):[page:0]

\begin{itemize}[leftmargin=1.5em]
\item \textbf{Level 0: Concrete agents.}
Multiple AI agents (LLM-based or otherwise) answer questions and produce structured replies.
Typical examples are role-based LLM agents such as a ``GRA optimist'' and a ``physics skeptic'', each configured with its own backend, system prompt and tool set.[page:0]

\item \textbf{Level 1: Forum orchestrator.}
The orchestrator manages multi-round debates, dispatches questions to agents, collects their replies and maintains a structured representation of the debate state, including extracted claims, embeddings and provenance.[page:0]

\item \textbf{Level 2: GRA core.}
The GRA core computes foam metrics \(\Phi\) on the debate state, aggregates them into a scalar objective \(J\), and proposes the next-round debate plan (e.g.\ whom to challenge, which contradictions to attack, where to request sources).[page:0]
\end{itemize}

The repository structure is aligned with this decomposition:[page:0]

\begin{verbatim}
gra-forum/
  gra_forum/
    core/          # GRA core and foam metrics
    orchestrator/  # Debate orchestrator
    agents/        # Base and concrete agent implementations
    models/        # Pydantic (or similar) data models
    utils/         # Embeddings and helper functions
  examples/        # Usage examples
  tests/           # Unit tests
  docs/            # Documentation
\end{verbatim}

\section{Foam Metrics and Global Objective}

The GRA core currently implements four foam components: conflict, vacuity, redundancy and noise.[page:0]

\paragraph{Conflict.}
\(\Phi_{\text{conflict}}\) measures unresolved contradictions between agents.
Operationally, GRA Forum represents each agent reply as a set of extracted claims plus sentence embeddings.
Pairs of claims with high semantic anti-correlation (e.g.\ via cosine distance in embedding space) that are not supported by sources or proofs contribute positively to \(\Phi_{\text{conflict}}\).[page:0]

\paragraph{Vacuity.}
\(\Phi_{\text{vacuity}}\) estimates the fraction of claims that lack verifiable support: missing citations, absence of quantitative checks, or unverifiable vague statements.
This encourages the orchestrator to push agents towards providing sources, calculations and explicit assumptions.[page:0]

\paragraph{Redundancy.}
\(\Phi_{\text{redundancy}}\) captures excessive repetition of similar content by different agents without added informational value.
High similarity of embeddings across claims from different agents, combined with low incremental novelty, increases this term.[page:0]

\paragraph{Noise.}
\(\Phi_{\text{noise}}\) represents lower-level errors such as invalid JSON, broken formatting or clear violations of basic constraints (e.g.\ impossible numerical values detectable by simple validators).
In the current MVP this component is implemented as a placeholder that can be refined over time.[page:0]

The global foam functional is then defined as
\begin{equation}
  J = w_c \,\Phi_{\text{conflict}} 
    + w_v \,\Phi_{\text{vacuity}}
    + w_r \,\Phi_{\text{redundancy}}
    + w_n \,\Phi_{\text{noise}},
\end{equation}
with user-configurable non-negative weights \(w_c, w_v, w_r, w_n\).[page:0]
The GRA principle is to design the orchestrator such that \(J\) tends to decrease over successive debate rounds, while still allowing non-zero constructive disagreement.[page:0]

\section{Orchestrator Logic}

The \texttt{DebateOrchestrator} class implements the high-level control loop:

\begin{enumerate}[leftmargin=1.5em]
\item \textbf{Initialization.}
The user selects a set of Level 0 agents and initializes the GRA core with weights for foam metrics and a maximum number of rounds.[page:0]

\item \textbf{Round execution.}
Given an initial question, the orchestrator dispatches it to all or a subset of agents, collects their replies, and updates the debate state (claims, embeddings, metadata).[page:0]

\item \textbf{Foam evaluation.}
The GRA core computes \(\Phi_{\text{conflict}}, \Phi_{\text{vacuity}}, \Phi_{\text{redundancy}}, \Phi_{\text{noise}}\) and \(J\) for the current state.
The history of \(J\) over rounds can be logged and visualized.[page:0]

\item \textbf{Planning the next round.}
Based on which foam component dominates, the core selects a strategy:
\begin{itemize}[leftmargin=1em]
\item if \(\Phi_{\text{conflict}}\) is high: \texttt{attack} --- confront agents holding opposing positions and request explicit resolution;
\item if \(\Phi_{\text{vacuity}}\) is high: \texttt{verify} --- ask agents to provide sources, calculations or formal arguments;
\item if \(\Phi_{\text{redundancy}}\) is high: \texttt{synthesize} --- ask one agent to synthesize and compress overlapping content, and redirect others to unexplored aspects.[page:0]
\end{itemize}

\item \textbf{Stopping.}
The debate stops when either the maximum number of rounds is reached or when \(J\) falls below a user-defined threshold, indicating that most parasitic foam has been removed.[page:0]
\end{enumerate}

A minimal usage example (from the \texttt{examples} directory) illustrates this loop:

\begin{verbatim}
from gra_forum.orchestrator import DebateOrchestrator
from gra_forum.agents import RoleBasedAgent
from gra_forum.core import GRACore

optimist = RoleBasedAgent(
    agent_id="gra_optimist",
    role="gra_optimist",
    llm_backend="openai",
    api_key="..."
)
skeptic = RoleBasedAgent(
    agent_id="physics_skeptic",
    role="physics_skeptic",
    llm_backend="anthropic",
    api_key="..."
)

core = GRACore(weights={
    "conflict": 0.40,
    "vacuity": 0.35,
    "redundancy": 0.20,
    "noise": 0.05,
})

orchestrator = DebateOrchestrator(
    agents=[optimist, skeptic],
    core=core,
    max_rounds=3,
)

question = "Can exponential growth in AI agents lead to " \
           "exponential growth in cognitive capacity?"
history = orchestrator.run_debate(question)
orchestrator.plot_J(save_path="debate_results.png")
\end{verbatim}[page:0]

\section{Extensibility and Physical Agents}

While the initial implementation focuses on text-based LLM agents, the design deliberately separates the GRA core from concrete agent backends.
Any agent that implements a simple interface
\begin{verbatim}
class BaseAgent:
    async def answer(self, question: str) -> AgentReply:
        ...
\end{verbatim}
can participate in the forum.
This opens the door to integrating ``physical AI'' agents, such as NVIDIA Omniverse or physics-informed neural network simulators, which can test the physical implications of debated hypotheses in digital twin environments.[web:242][web:247][web:249]

In such settings, foam metrics can be augmented by physical inconsistency terms, e.g.\ deviation between simulated dynamics and basic conservation laws, or disagreement between agents' verbal claims and numerical simulation results.[web:247][web:250]

\section{Conclusion and Future Work}

GRA Forum is a first step towards treating multi-agent AI debates as controlled dynamical systems rather than ad-hoc chains of prompts.
By explicitly defining and minimizing argumentative foam, the framework aims to:
(i) preserve productive disagreement,
(ii) suppress non-productive noise and redundancy,
and (iii) provide quantitative traces of debate quality over time.[page:0]

Future work includes:
adding richer claim extraction and contradiction detection,
integrating external verifiers (symbolic math, code execution, physical simulators),
and exploring learning-based GRA cores that adapt weights and strategies from data.[web:242][web:247]

\section*{Acknowledgements}

The author thanks the broader AI and open-source communities for inspiration and tools that make such orchestration experiments possible.

\bibliographystyle{plain}
\begin{thebibliography}{9}

\bibitem{gra_forum}
AAA AAA.
\newblock GRA Forum: AI Agent Debate Orchestrator.
\newblock GitHub repository, \url{https://github.com/qqewq/gra-forum}, 2024.

\bibitem{zenodo}
AAA AAA.
\newblock GRA Forum: AI Agent Debate Orchestrator.
\newblock Zenodo, \url{https://doi.org/10.5281/zenodo.19158686}, 2024.

\end{thebibliography}

\end{document}
